\documentclass[12pt]{article}
\usepackage[utf8]{inputenc}
\usepackage[T1]{fontenc}
\usepackage{hyperref}
\title{Yavitero Language Resources for K-12 Curriculum Development}
\author{Compiled by Claude Sonnet 5.5 (claude-sonnet-5-5) for Andr\'as Kornai}
\date{}
\begin{document}
\maketitle

\section{Introduction}
Yavitero (ISO 639-3 yvt) is an extinct Maipurean (Arawakan) language of the Atabapo River basin, Venezuela. Sources say it was practically extinct, with one known elderly speaker reported in 2000, and they disagree on the end date. Only reference-level sources were found.

\section{Primary Educational Resources}
None found. No school or revitalization programme appeared in the searches.

\section{Linguistic and Reference Materials}
\begin{itemize}
\item Wikipedia entry on Yavitero: \url{https://en.wikipedia.org/wiki/Yavitero_language}
\item Berkeley SAPHON phonological inventory for Yavitero, which cites the Mosonyi works: \url{https://linguistics.berkeley.edu/saphon/en/inv/Yavitero.html}
\item Endangered Languages Project entry, with a bibliography: \url{https://endangeredlanguages.com/lang/1640} and \url{https://endangeredlanguages.com/lang/1640/bibliography}
\item The main sources named in the search summaries are J. Mosonyi's 1987 dissertation \emph{El idioma yavitero: ensayo de gramática y diccionario} (Universidad Central de Venezuela) and the Yavitero chapter by Mosonyi, Mosonyi and Largo in \emph{Manual de lenguas indígenas de Venezuela} (2000). No online copy of either appeared in the search results.
\item Related Avane language, known only from a short word list: \url{https://en.wikipedia.org/wiki/Avane_language}
\end{itemize}

\section{Religious and Community Resources}
None found.

\section{Digital and Audio Resources}
None found.

\section{Curriculum Development Context}
With no living community of speakers reported, any heritage curriculum would have to be reconstructive. It would rest on the Mosonyi grammar and dictionary, which would need to be located in print, and on related Arawakan languages.

\end{document}